% TOP_PROBLEM: {"id": 3366, "title": "The Riemann Hypothesis", "category_id": 1, "category": {"id": 1, "name": "number_theory", "display_name": "Number Theory", "description": "Properties of integers, prime numbers, Diophantine equations.", "slug": "number-theory", "order_index": 1, "created_at": "2026-07-31T15:26:25.670Z"}, "status": "open", "problem_number": "OPG-573", "set_id": 12}
% FIRST_POSTED: 2026-10-01
% ENTRY_KIND: statement_only
% UnsolvedMath ID 3366; source catalogue code OPG-573
% !TeX program = lualatex
% Definition and sources only; this file is not an LLM solution attempt.
\documentclass[11pt,a4paper]{article}
\usepackage[margin=25mm]{geometry}
\usepackage{fontspec}
\setmainfont{lmroman10-regular.otf}[BoldFont=lmroman10-bold.otf,
 ItalicFont=lmroman10-italic.otf,BoldItalicFont=lmroman10-bolditalic.otf]
\usepackage{amsmath,amssymb,mathtools}
\usepackage{unicode-math}
\setmathfont{latinmodern-math.otf}
\AtBeginDocument{\let\setminus\smallsetminus}
\usepackage{xurl,hyperref}
\hypersetup{colorlinks=true,urlcolor=blue}
\setcounter{secnumdepth}{0}
\setlength{\parindent}{0pt}
\setlength{\parskip}{5pt}
\setlength{\emergencystretch}{3em}
\begin{document}
\section*{The Riemann Hypothesis}\label{3366}
{\small UnsolvedMath ID 3366 · OPG-573 · Number Theory · OpenGarden}
\subsection{Definitions and mathematical statement}
For $s\in\mathbb C$ with $\operatorname{Re}s>1$, define
\[
 \zeta(s)=\sum_{n=1}^{\infty}n^{-s},
 \qquad n^{-s}=\exp(-s\log n),
\]
using the real natural logarithm of the positive integer $n$. The
Riemann zeta function on other arguments means the unique meromorphic
continuation of this function to $\mathbb C$. It has a simple pole at
$s=1$. The critical strip is $0<\operatorname{Re}s<1$, and the critical
line is $\operatorname{Re}s=1/2$.

The Riemann hypothesis asserts
\[
 \zeta(\rho)=0,\quad 0<\operatorname{Re}\rho<1
       \quad\Longrightarrow\quad \operatorname{Re}\rho=\tfrac12.
\]
The quantifier ranges over all complex zeros in the strip, regardless
of their imaginary height or multiplicity. No assertion that the zeros
are simple is included. The familiar zeros at negative even integers
are outside this strip and are not counterexamples. Likewise the pole
at one is not a zero. Evaluating the original Dirichlet series inside
the strip without continuation would not define the function used here.
The target concerns this single zeta function; corresponding assertions
for Dirichlet, Dedekind, or other $L$-functions have wider scope.
\subsection{Short English statement}
Must every zero of the Riemann zeta function inside its critical strip lie on the line with real part one half?
\subsection{Sources}
\begin{itemize}
\item[S1] ULAM.AI and the UnsolvedMath Contributors, supplied problems.json, record 3366 (OPG-573), OpenGarden collection. Catalogue identity and source transcription; CC BY 4.0.
\url{https://huggingface.co/datasets/ulamai/UnsolvedMath}.
\item[S2] Open Problem Garden, The Riemann Hypothesis. Original problem and discussion; the mathematical formulation below is expanded from the supplied source record.
\url{http://www.openproblemgarden.org/op/the_riemann_hypothesis}.
\end{itemize}
\end{document}
